\chapter{Stato dell'arte}
\label{ch:stato_arte}

% BOZZA — struttura proposta, da validare con la relatrice.

\section{Generazione automatica di test: approcci classici}
% TODO: esecuzione simbolica, search-based testing (EvoSuite), ecc.
% Qui si colloca anche Klara: generazione di scheletri di test tramite il
% solver SMT Z3, capace di trovare input che coprono ogni ramo (boundary
% value trovati automaticamente), ma limitata a un sottoinsieme di Python.

\section{Generazione di test con modelli linguistici}
% TODO: panoramica LLM per il testing.

\section{Il benchmark TestGenEval}
% Punti chiave dal paper (ICLR 2025):
% - costruito su SWE-bench: 68.647 test da 1.210 coppie file codice-test,
%   11 repository Python reali e mantenuti;
% - due task: generazione dell'intera test suite e test completion
%   (primo/ultimo/test aggiuntivo);
% - metriche: pass@1/pass@5, coverage, mutation score (bug sintetici
%   iniettati nel codice: quanti ne scoprono i test?);
% - risultato principale: i modelli faticano sui progetti reali; il
%   migliore (GPT-4o) si ferma al 35,2% di coverage media e 18,8% di
%   mutation score;
% - causa principale: difficoltà a ragionare sull'esecuzione, errori nelle
%   assert sui cammini complessi;
% - il test completion è più facile, ma i modelli aggiungono meno dell'1%
%   di coverage sull'ultimo test.

\section{Limiti degli approcci esistenti}
% TODO: i modelli generano "alla cieca", in un colpo solo, senza osservare
% l'esecuzione: è il limite che la tesi affronta.
